\section{Clasidicador Random Forest}
En este trabajo se decidió usar el algoritmo de Random Forest el cual es un algoritmo de
aprendizaje automático (machine learning) utilizado para clasificar y pronosticar. Esta técnica
utiliza el concepto de árboles de decisión, construyendo una colección de árboles de decisión y
agregando sus resultados para generar la predicción definitiva. Cada árbol de decisión dentro de un
bosque aleatorio se construye utilizando subconjuntos aleatorios de datos, y cada árbol individual
5
se entrena en una parte de todo el conjunto de datos. Posteriormente, los resultados de todos
los árboles de decisión se amalgaman para derivar el pronóstico final. Uno de los beneficios de
los bosques aleatorios es su capacidad para manejar datos desequilibrados. Además, los bosques
aleatorios mitigan la cuestión del sobre ajuste mediante la capacitación de varios árboles de decisión
en subconjuntos aleatorios de datos, lo que mejora su capacidad para generalizar a datos novedosos;
una de las aplicaciones de este método es para diagnosticar pacientes y predecir resultados de
pacientes salman2024random.
Ademas se complementara con el método Leave One Out el cual es una técnica de estimación
estadística que se aplica en diferentes campos que requieran una evaluación del rendimiento de un
algoritmo de aprendizaje. Este se trata de un caso especial del método de validación cruzada k-fold
que implica la partición del conjunto de datos originales k subconjuntos de tamaño (aproximada-
mente)igual. El modelo se entrena k veces, utilizando cada subconjunto por turno como conjunto de
prueba, y los subconjuntos restantes como conjunto de entrenamiento. La precisión total se puede
obtener promediando los valores de precisión calculados en cada subconjunto. En el caso de leave
one out este es una validación cruzada k-fold calculada con k = n, donde n es el tamaño del con-
junto de datos original. Por lo tanto cada conjunto de prueba tiene tamaño 1, lo que implica que el
modelo se entrena n veces brovelli2008accuracy.
